\section{Discussion and Application Development}
The architecture connects object retrieval through explicit subsystem evidence: vision proposes a candidate, task authority permits handling, navigation establishes a pickup region, the base confirms rest, and manipulation confirms grip and release. The mission model exposes these obligations where hardware feedback remains to be developed. Its nominal trace demonstrates logical progress under stated assumptions, while source analysis identifies what the actuator gateway can express.

The next application step is to instantiate these observations with measured quantities. This requires a task-specific visual dataset and held-out recognition results; calibrated camera--LiDAR--arm frames and contact geometry; feedback for base rest, grip, and release; and timed physical runs covering navigation, pickup, transport, and deposit. Sensor-age limits must follow measured acquisition and transport behavior. Force, payload, containment, and environmental compatibility must be evaluated for the designated object class. A general RGB category detector cannot establish an unknown material's hazardous properties.

The results establish source and model properties; physical obstacle avoidance and hazardous-environment suitability require application trials. Nonhazardous surrogates support the initial benchmark, followed by object- and environment-specific qualification as illustrated by prior systems~\cite{bird2022vega,kemp2022stretch}.
